\SourceLeaf{356}%
small that for most measurements one will hardly consider them worth even the slight trouble of taking them into account.

The memoir fully develops the theory of all these calculations and others related to them, illustrates it with a worked example from start to finish, and supplies an auxiliary table. The table is intended in the first instance for the zone containing the Hanoverian triangulation network, but it can also be used directly for surveys that extend far beyond that zone. It spans twelve degrees of latitude, with the chosen standard parallel at \(52^0\,40'\) latitude in the middle. The table is calculated to an accuracy sufficient even for a triangulation network to be calculated with ten-figure logarithms, and thus to much greater accuracy than will be retained in most cases. Still, the slight saving in space that would have resulted from dropping a couple of decimal places seemed too small to justify leaving any out in print.

What is perceptible and what is imperceptible in calculations are relative notions. It is therefore worth putting a definite measure on them with a couple of examples taken from the memoir.

The largest triangle in the Hanoverian triangulation network, which is also the farthest from the standard parallel, is the triangle formed by Brocken, Hohenhagen and Inselsberg. It therefore also contains the largest values of the direction reductions. They occur on the Hohenhagen-Inselsberg side, where the azimuth reduction is \(-0''00332\) at the first endpoint and \(+0''00428\) at the other. In the entire network there are only two other triangle sides for which the reductions exceed \(0''001\). On all the remaining sides they stay below that amount.

The largest principal triangle in the trigonometrical surveys of Switzerland is the one between Chasseral, Suchet and Berra. It just reaches the southern limit of the auxiliary table, so the direction reductions can still be calculated with it. The largest reduction is the one that applies to the azimuth of Chasseral at Suchet, and it amounts to \(+0''06221\).

It follows that throughout the zone containing the Hanoverian triangulation network, the reduction can be left out entirely when the calculation is carried to
